\section*{Chapitre \ref{ch:pl:observability-and-incident-response} --- Observabilité et réponse aux incidents}

\begin{solution}{exo:pl:observability-and-incident-response:1}
$0.001 \times 30 \times 23\,400 = 702$ secondes non conformes~; le mois calme en a consommé 115, soit 16.4~\% du budget, réparties en 18 pauses de plus de cinq secondes.
\end{solution}

\begin{solution}{exo:pl:observability-and-incident-response:2}
Parce qu'aucun processus ne s'arrête~: pendant le blocage, les processus attendent des données qui n'arrivent pas, lors des blocages intermittents ils attendent vingt secondes par
minute, et, avec le consommateur lent, le service de risque est occupé. Un heartbeat mesure le processus, et non le service que reçoivent les utilisateurs.
\end{solution}

\begin{solution}{exo:pl:observability-and-incident-response:3}
56 secondes~: les cinq premières secondes restent conformes (données âgées de cinq secondes au plus)~; la fenêtre d'une heure doit alors contenir 1.44~\% de 3\,600 secondes non conformes,
soit 51.8, donc l'alerte arrive à la 52\up{e} seconde non conforme, la fenêtre de cinq minutes étant depuis longtemps révolue.
\end{solution}

\begin{solution}{exo:pl:observability-and-incident-response:4}
Chaque blocage dure vingt secondes, puis les données se rafraîchissent, si bien que leur âge ne dépasse jamais vingt secondes. Un tiers des secondes sont non conformes, ce qui
consomme vite le budget, mais aucune seconde isolée ne franchit trente.
\end{solution}

\begin{solution}{exo:pl:observability-and-incident-response:5}
Elle exige que le problème soit en cours au moment de l'évaluation~: sans cette condition, une fenêtre d'une heure continue de déclencher des alertes pendant une heure après la fin d'un incident, et
une salve survenue tôt dans la fenêtre continue d'appeler après le rétablissement. La fenêtre courte fait cesser l'alerte peu après le rétablissement du service.
\end{solution}

\begin{solution}{exo:pl:observability-and-incident-response:6}
21 minutes, de la dernière cotation à 12~h~03 jusqu'aux données fraîches à 12~h~24, mesurées sur le service que reçoit le système de risque. Un temps moyen de
réparation (livre~14, chapitre~15) désigne la moyenne, pour un composant, du délai entre la panne et la réparation sur de nombreuses pannes~; il s'agit ici d'un seul incident sur un seul service.
\end{solution}

\begin{solution}{exo:pl:observability-and-incident-response:7}
Le seuil de cinq secondes double et atteint 36 fausses alertes par mois, celui de trente secondes triple et atteint 3~; les règles de vitesse de consommation restent à une, car les
pauses restent courtes et ne consomment chacune qu'un peu de budget.
\end{solution}

\begin{solution}{exo:pl:observability-and-incident-response:8}
Les services étaient sains et inactifs~: ils se déclaraient eux-mêmes en bon état, et non sur ce qu'ils livraient. Les contrôles de santé doivent être complétés par des SLI issus de
du côté des utilisateurs --- ici l'âge des données --- faute de quoi une panne en amont ressemble à une journée calme.
\end{solution}

\begin{solution}{pb:pl:observability-and-incident-response:1}
\begin{enumerate}
\item \omterm{def:pl:observability-and-incident-response:metric}{Métriques}~: indiquent si quelque chose ne va pas, partout et à faible coût. Journaux~: ce qui s'est passé dans un service. Traces~: où le temps a été consommé à travers les services pour une
unité de travail.
\item Fractions de milliseconde à la capture, au stockage et aux positions, puis dans la file d'attente du risque~: 50~ms en temps normal, 60~s pendant le consommateur lent~; le calcul de
risque ne prend que 1~ms.
\item Parce qu'un service peut être vivant et inutile~: les utilisateurs ont besoin de données fraîches, et non de processus actifs.
\item 702 secondes non conformes par mois~; les jours calmes en consomment 115.
\item Un SLO est la cible propre à la société, fixée plus strictement~; un SLA est un engagement envers un client, assorti de pénalités.
\item Un blocage du flux à partir de 12~h~03 pendant 21 minutes, des blocages de vingt secondes chaque minute à partir de 11~h~48 pendant 15 minutes, un consommateur lent à partir de 14~h~00.
\item Cinq secondes~: les trois, en 5, 5 et 100~s~; trente secondes~: le blocage en 30~s et le consommateur lent en 600~s, mais pas les blocages intermittents~;
soixante secondes~: le blocage en 60~s, le consommateur lent en 1\,200~s, mais pas les blocages intermittents.
\item 18, 1 et 0 pour les trois seuils, aucune pour le heartbeat, une pour les règles de vitesse de consommation.
\item Les trois, en 56, 191 et 151 secondes.
\item C'est la pause de 56 secondes, qui consomme une part réelle du budget et mérite un examen.
\item Tient à jour l'état de l'incident et attribue les rôles~; ne fait pas le travail technique.
\item Exposition~: la société négocie-t-elle sur des données périmées, et faut-il arrêter les stratégies --- des kill switches et des limites de risque, avant tout diagnostic.
\item 12:03:00 dernière cotation~; 12:03:05, 12:03:30, 12:03:56 et 12:04:00 les quatre alertes~; 12:24:00 le flux reprend et les données sont fraîches.
\item 56 secondes de \omterm{def:pl:observability-and-incident-response:ttd}{délai de détection} avec la règle de vitesse de consommation, 21 minutes de \omterm{def:pl:observability-and-incident-response:ttd}{délai de rétablissement}.
\item Une chronologie, les causes contributives sans recherche de coupable, et les modifications apportées au système.
\item \textbf{Résultat nommé.} Un blocage complet du flux est détecté en 5, 30 et 60 secondes par des seuils statiques à ces âges, et en 56 secondes par
les règles de vitesse de consommation~; un consommateur lent en 100, 600 et 1\,200 secondes contre 151~; les blocages intermittents uniquement par le seuil de cinq secondes (5
secondes) et les règles de vitesse de consommation (191). Sur trente jours calmes, les seuils envoient 18, 1 et 0 fausses alertes par mois, les règles de vitesse de consommation une~;
un heartbeat de processus ne détecte rien.
\item Les règles de vitesse de consommation pour les alertes, avec un ticket sur une vitesse de consommation de 1 sur trois jours pour les dégradations lentes.
\item Un SLI portant sur l'âge des données utilisées par chaque consommateur.
\item Des SLI d'âge des données partout, des alertes déclenchées selon la vitesse de consommation, un kill switch sur les données périmées, et le rejeu conservé comme test de l'alerte.
\item Qu'une personne doit agir maintenant, sur quelque chose que les utilisateurs du service peuvent ressentir.
\end{enumerate}
\end{solution}

\begin{solution}{iq:pl:observability-and-incident-response:1}
Un contrôle de processus vérifie si un programme tourne~; un contrôle de service vérifie si ses utilisateurs obtiennent ce dont ils ont besoin, mesuré de leur côté (âge des données,
taux d'erreur, latence). Seul le second voit une panne en amont ou un ralentissement silencieux.
\iqlookfor{Indicateurs côté utilisateur.}
\end{solution}

\begin{solution}{iq:pl:observability-and-incident-response:2}
Un SLI mesure le service (la part de secondes où les données de marché ont cinq secondes d'âge au plus)~; le SLO en est la cible (99.9~\% sur 30 jours
de négociation)~; le \omterm{def:pl:observability-and-incident-response:budget}{budget d'erreur} correspond à ce que la cible permet (702 secondes non conformes par mois), consommé par les incidents et les changements risqués.
\iqlookfor{Un indicateur concret, sa cible et son budget.}
\end{solution}

\begin{solution}{iq:pl:observability-and-incident-response:3}
Supprimer ou déclasser toute alerte qui ne nécessitait pas l'intervention immédiate d'une personne~; déclencher les alertes selon la vitesse de consommation du SLO plutôt que sur des seuils bruts~; associer une \omterm{def:pl:observability-and-incident-response:oncall}{procédure d'exploitation} (runbook) à chaque page~;
revoir les alertes chaque semaine.
\iqlookfor{Une alerte implique une action.}
\end{solution}

\begin{solution}{iq:pl:observability-and-incident-response:4}
Il faut la tracer~: propager un identifiant de trace avec le message dans tous les services, enregistrer un \omterm{def:pl:observability-and-incident-response:trace}{span} par étape et lire l'arbre~; une attente en file
apparaît comme un seul \omterm{def:pl:observability-and-incident-response:trace}{span} long.
\iqlookfor{\omterm{def:pl:observability-and-incident-response:trace}{Traçage distribué}.}
\end{solution}

\begin{solution}{iq:pl:observability-and-incident-response:5}
SLI et SLO côté utilisateur par service (âge des données, complétude, latence)~; \omterm{def:pl:observability-and-incident-response:metric}{métriques}, \omterm{def:pl:observability-and-incident-response:metric}{journalisation structurée} et traces avec identifiants propagés~;
\omterm{def:pl:observability-and-incident-response:burn}{alertes de vitesse de consommation} et tickets de consommation lente~; \omterm{def:pl:observability-and-incident-response:oncall}{rotations d'astreinte} avec des \omterm{def:pl:observability-and-incident-response:oncall}{procédures d'exploitation} qui commencent par l'exposition~; désignation d'un \omterm{def:pl:observability-and-incident-response:ic}{responsable d'incident}~; \omterm{def:pl:observability-and-incident-response:blameless}{revues sans recherche de coupable} avec
actions suivies~; incidents rejoués comme tests.
\iqlookfor{SLO, vitesses de consommation, pratique de l'incident.}
\end{solution}
